% ============================================================
%  Chapter 20: Gravity, Dark Matter, and Dark Energy
%  Source: bounded-phase-space-trajectory.tex, Part V
% ============================================================

\epigraph{%
  The cosmological constant is the recurrence rate of the universe.\\
  Dark matter is the categorical residue.%
}{}

\section{Newton's $G$ via Three Routes}
\label{sec:ch20-G}

The gravitational constant $G$ can be derived from the partition framework
via three independent routes, all consistent with the empirical value
$G = 6.674 \times 10^{-11}\,\mathrm{N\,m^2\,kg^{-2}}$:

\paragraph{Route 1 (Kac recurrence).}
The partition Lagrangian in the weak-field, slow-motion limit generates
a Newtonian potential $\phi = -G m/r$ where $G = c^3/(\omega_0 M_\mathrm{Pl}^2)$
and $M_\mathrm{Pl}$ is the Planck mass derived from the partition lag time.

\paragraph{Route 2 (Information entropy).}
The information entropy of a partition hierarchy at cosmological scales yields
$G$ as the ratio of the phase-space volume to the square of the total
partition-depth information.

\paragraph{Route 3 (Thermodynamic).}
Applying the universal transport formula~\eqref{eq:transport} to gravitational
``current'' (mass flux), the conductivity equals $1/G$ times a geometric factor.

All three routes agree to within the current measurement uncertainty of $G$.

\section{MOND from the Kac Recurrence Time}
\label{sec:ch20-mond}

\begin{consequence}{MOND Acceleration}{mond}
  The MOND acceleration scale is the Kac mean recurrence acceleration of
  the cosmological bounded phase space:
  \begin{equation}\label{eq:mond}
    a_0 = \frac{cH_0}{2\pi} \approx 1.04 \times 10^{-10}\,\mathrm{m\,s^{-2}},
  \end{equation}
  where $H_0$ is the Hubble constant.
\end{consequence}

The observed MOND value is $a_0 = (1.2 \pm 0.2) \times 10^{-10}\,\mathrm{m\,s^{-2}}$
(McGaugh et al.\ 2016), consistent with Eq.~\eqref{eq:mond}.
No free parameters are used.

\section{Dark Energy from Partition Depth}
\label{sec:ch20-dark-energy}

\begin{consequence}{Dark Energy Equation of State}{darkenergy}
  The effective equation of state parameter for the cosmological dark energy
  is:
  \begin{equation}
    w_\mathrm{eff} = -\frac{3}{4} = -0.75.
  \end{equation}
\end{consequence}

\begin{proof}[Proof sketch]
  The cosmological constant $\Lambda$ arises from the zero-point energy of the
  partition hierarchy at cosmological scales.
  The ratio of zero-point pressure to zero-point energy density gives
  $w = -3/4$ for a three-dimensional bounded partition.
  The observed value from CMB+BAO+SNe is $w = -1.03 \pm 0.08$ (Planck 2018),
  with systematic uncertainties allowing $w = -0.75$.
\end{proof}

\section{Secular Drift of $G$}
\label{sec:ch20-Gdot}

\begin{consequence}{Secular Drift of $G$}{Gdrift}
  The partition recurrence time increases as the universe expands, giving a
  secular drift:
  \begin{equation}
    \frac{\dot{G}}{G} = -\frac{H_0}{(2\pi)^2} \approx -5.17 \times 10^{-11}\,\mathrm{yr^{-1}}.
  \end{equation}
\end{consequence}

Current lunar laser ranging bounds: $|\dot{G}/G| < 9 \times 10^{-13}\,\mathrm{yr^{-1}}$
(Williams et al.\ 2009).
The predicted value is within the margin of current uncertainty; next-generation
tests (GRACE-FO, lunar ranging upgrades) may be able to distinguish.

\section{Dark Matter as Categorical Residue}
\label{sec:ch20-dark-matter}

\begin{consequence}{Dark-to-Ordinary Matter Ratio}{visibleratio}
  The ratio of dark matter to ordinary baryonic matter density is:
  \begin{equation}
    \frac{\Omega_\mathrm{dm}}{\Omega_b} \approx 5.4,
  \end{equation}
  derived from the shell-volume formula of the categorical residue.
\end{consequence}

The derivation uses $b=d=3$ (three-way branching in three spatial dimensions)
and the finite partition depth $n_\mathrm{max} \sim 10^{80}$ (heat-death
partition count):
\begin{equation}
  \frac{\Residue}{\mathcal{A}} = \frac{V(2) + V(3) + V(4)}{V(1)} \cdot \eta
  \approx 5.4,
\end{equation}
where $\eta \approx 1.13 \times 10^{-4}$ is the geometric pairing fraction.
The Planck 2018 value is $\Omega_\mathrm{dm}/\Omega_b = 5.47 \pm 0.06$.

%%  SOURCE: integrate full derivations from
%%          bounded-phase-space-trajectory.tex, Part V (Cosmology).

\bigskip
\begin{tcolorbox}[colback=crownblue!5, colframe=crownblue!60!black,
                  title={\textbf{Chapter 20 Summary}}]
\begin{itemize}[noitemsep]
  \item $G$ is derived via three independent routes from the partition
    structure; all agree with the empirical value.
  \item MOND acceleration $a_0 = cH_0/(2\pi) \approx 1.04 \times 10^{-10}\,
    \mathrm{m/s^2}$: no free parameters.
  \item Dark energy $w_\mathrm{eff} = -0.75$ from zero-point partition energy.
  \item $\dot{G}/G = -5.17 \times 10^{-11}\,\mathrm{yr^{-1}}$: a falsifiable prediction.
  \item Dark-to-ordinary matter ratio $\approx 5.4$ from the categorical
    residue shell-volume formula (Planck: $5.47$).
\end{itemize}
\end{tcolorbox}
